% Target ~350 words. Replaces preliminary stub.

\subsection{Geometry as clinical signal}
The headline finding is that severity correlates with a purely
\emph{geometric} property of frozen WavLM representations: the
length-normalised DTW distance from a control-built reference
trajectory. No fine-tuning, no labelled training, no classifier.
The implicit articulatory manifold learned during 94k hours of
neurotypical pre-training is sufficient to expose dysarthric
deviation as off-manifold travel. Replacing DTW with mean-pool L2
collapses $\rho$ by 40\%, and replacing WavLM with MFCC trajectories
yields a higher in-corpus global $\rho$ but a lower dysarthric-subgroup
$\rho$ ($0.282$ vs $0.394$; both $n = 8$, non-significant, read as
directional only). These two ablations bracket the contribution of
\emph{trajectory structure} and of \emph{pre-trained SSL geometry}
respectively.

A trained linear probe on mean-pooled WavLM attains comparable accuracy
(LOSO $\rho = 0.760$--$0.775$, RMSE $0.581$ vs.\ ATD's $0.756$/$0.671$;
Table~\ref{tab:main}); ATD's contribution is therefore \emph{not} higher
accuracy but label-free zero-shot operation, per-phoneme
interpretability, and real-time deployment.

\subsection{Layer choice and the falsified specificity hypothesis}
Layer~9 is the primary layer because it is robust to early-layer
speaker-identity leakage: under leave-one-control-out references the
Layer-9 headline is unchanged ($0.756 \to 0.756$) whereas Layer~1
inflates ($0.775 \to 0.687$). The pre-registered controlled-perturbation
study does \emph{not} support an articulatory-specificity claim at either
layer: pink noise outranks formant shift at Layer~9
($0.77$ vs $0.17$) \emph{and} at Layer~1 ($0.83$ vs $0.40$), which is
exactly the locked falsification condition (criterion~\#2). We therefore
report Track~C as a falsified hypothesis (Limitations) and do not use the
perturbation profile to recommend a layer. The one robust positive from
the perturbation sweep is that the trajectory metric detects partial
muting (frame dropout) far better than pooled MFCC
($\Delta\rho{=}+0.25$), i.e.\ it uses temporal structure MFCC discards ---
but that is a statement about trajectories, not about formant-level
articulatory selectivity. A separate, corpus-specific observation on IPVS
(XLS-R-300m) is that the within-corpus PD-vs-HC signal there is carried by
the acoustic-near Layer~1 rather than Layer~9; we flag this as
backbone- and corpus-specific, subject to the same early-layer leakage
caveat, and not a general layer recommendation. The headline TORGO number
is reported at Layer~9 ($\rho{=}0.756$), the leakage-robust primary
layer.

\subsection{Trajectory-length regime}
Direct cross-corpus transfer to UASpeech fails ($\rho{=}{-}0.22$,
$p{=}0.43$) and we report this as a falsified hypothesis. The
trajectory-length analysis (TORGO mean 155 frames vs UASpeech mean
127 frames at 50~fps) provides a mechanistic, falsifiable
explanation: ATD requires a minimum temporal regime for DTW signal
to stabilise. Stitched-sentence UASpeech experiments (in progress)
will test this directly.

\subsection{F04 as a clinical finding, not an anomaly}
Removing F04 shifts $\rho$ from $0.756$ to $0.855$. F04 has
documented preserved intelligibility despite a categorical
``severe'' label (spastic dysarthria with intact phonatory
control). ATD ranks her near controls, agreeing with the perceptual
record but disagreeing with the ordinal label. We argue this is the
substantive contribution of execution-quality metrics: they
disambiguate cases where the diagnostic category over-predicts
intelligibility loss.

\subsection{Cross-disorder and clinical anchoring}
We extend the analysis to Italian Parkinsonian speech (IPVS,
$n{=}24$ PD + $n{=}37$ HC) using XLS-R-300m features and per-corpus
CMVN with a TORGO-fit shared PCA basis. The within-corpus binary
discrimination (PD vs HC) provides a falsifiable test of the
cross-disorder claim. Partial Spearman of ATD vs UASpeech 3-listener
intelligibility, controlling for word identity, anchors ATD to a
clinical perceptual instrument (Analysis C1).

\subsection{Future work}
Three threads. First, a real-time streaming implementation
(sub-200~ms latency, online subsequence DTW) for clinic deployment.
Second, phoneme-level decomposition of the DTW path for per-token
SLP feedback. Third, a longitudinal pilot ($n{=}20$ PD, 5 visits
over 12 months) to validate within-subject sensitivity to therapy
and medication state.
